\section{Conclusion}
\label{sec:conclusion}

We asked whether a model can tell when it lacks the information to act safely, and choose among
acting, asking, refusing, and deferring. To answer it we built \carb, an 840-item four-way
action-routing benchmark with source-inherited labels, and held the items fixed while varying only
how the decision is elicited, across five regimes and five models.

The answer is a qualified yes with a specific correction. Models can be prompted into a usable
four-way policy --- 0.598--0.721 accuracy against 0.231 random, with over-commitment falling from
45--58\% to 9--25\% --- but not by the mechanism the hypothesis names. A verbalised confidence,
given the best router that can be fitted to it, loses to typed prompting by 0.158--0.275 on every
model and is at chance at telling apart the three reasons for not acting. The recognition that
fails is specific: models judge safety at 0.893 and their own capability at 0.943, and information
sufficiency at 0.565. Announced stakes move the decision criterion and not the discrimination, with
an ambiguity $\times$ stakes interaction of exactly zero. Fine-tuning an open 4B model reaches 0.794
in-distribution and 0.470 out of it, where a frozen linear probe on the untrained model already
reaches 0.779 --- so training mostly supplied a read-out for a signal the representation already
carried.

{\bf The takeaway:} the bottleneck in deciding whether to ask is structural, not statistical. Give
a model the action space and it routes; give it a confidence scale and it cannot, because three
reasons for not acting do not fit on one axis. And until a routing policy beats \emph{always asking}
at a realistic error cost, the honest baseline for any such system is the trivial one.

{\bf Future work,} in priority order. (1) A stakes manipulation with \emph{experienced}
consequences --- an executable task where a wrong action actually fails a test --- since announced
stakes are the weakest link here. (2) Training directly on the ASK-F1 objective rather than the
typed label, following the transfer result of \citet{trinh2026hilbench}. (3) Probing what the frozen
representation decodes that does not survive a source change, which is the crux of the negative
transfer result. (4) Multi-turn: whether to ask is not the same decision as when to stop.
